% TOP_PROBLEM: {"id": 1409, "title": "Universal Point Sets", "category_id": 3, "category": {"id": 3, "name": "graph_theory", "display_name": "Graph Theory", "description": "Problems involving graphs, networks, and their properties.", "slug": "graph-theory", "order_index": 3, "created_at": "2026-07-31T15:26:25.671Z"}, "status": "open"}
% FIRST_POSTED: null
% ENTRY_KIND: statement_only
% Imported from: unsolvedmath_additional_500.tex; UnsolvedMath ID 1409
% !TeX program = lualatex
% Compile twice: lualatex 1409.tex
% Self-contained: no external bibliography, images, or \input files.
% Numeric section labels and visible numbers are original UnsolvedMath IDs.
\documentclass[11pt,a4paper]{article}
\usepackage[margin=24mm,headheight=14pt]{geometry}
\usepackage{fontspec}
\setmainfont{Latin Modern Roman}
\setsansfont{Latin Modern Sans}
\setmonofont{Latin Modern Mono}
\usepackage{amsmath,amssymb,mathtools}
\usepackage{unicode-math}
\setmathfont{Latin Modern Math}
\usepackage{microtype}
\usepackage{enumitem,xurl,needspace}
\usepackage[dvipsnames]{xcolor}
\usepackage{hyperref}
\hypersetup{unicode=true,colorlinks=true,linkcolor=MidnightBlue,urlcolor=MidnightBlue,
 pdftitle={UnsolvedMath Problem 1409},pdfauthor={Source-annotated compilation},bookmarksnumbered=true}
\usepackage{bookmark,tocloft,titlesec,fancyhdr}
\setlength{\cftsecnumwidth}{4.2em}
\cftsetpnumwidth{2.5em}
\cftsetrmarg{4em}
\titleformat{\section}{\Large\bfseries\raggedright}{\thesection}{0.7em}{}
\pagestyle{fancy}\fancyhf{}
\fancyhead[L]{\small\sffamily Additional UnsolvedMath statements}
\fancyhead[R]{\small Original catalogue IDs}
\fancyfoot[C]{\thepage}
\setlength{\parindent}{0pt}
\setlength{\parskip}{5pt plus 1pt}
\setlength{\emergencystretch}{3em}
\setcounter{tocdepth}{1}\setcounter{secnumdepth}{1}
\setlist[itemize]{leftmargin=3em,itemsep=4pt,topsep=4pt,parsep=0pt}
\allowdisplaybreaks
\newcommand{\N}{\mathbb{N}}
\newcommand{\Z}{\mathbb{Z}}
\newcommand{\Q}{\mathbb{Q}}
\newcommand{\R}{\mathbb{R}}
\newcommand{\C}{\mathbb{C}}
\newcommand{\F}{\mathbb{F}}
\newcommand{\A}{\mathbb{A}}
\newcommand{\PP}{\mathbb{P}}
\newcommand{\E}{\mathbb{E}}
\newcommand{\eps}{\varepsilon}
\newcommand{\dd}{\,\mathrm d}
\newcommand{\sref}[2]{\hyperlink{source:#1:#2}{[S#2]}}
\newif\ifshowcatalogue
\showcataloguetrue
\newcommand{\cataloguescope}[1]{\ifshowcatalogue
 \paragraph{Statement recorded in the supplied source.}
 #1\fi}
\begin{document}
% UnsolvedMath ID 1409; source catalogue code GRAPH-022

\setcounter{section}{1408}

\section{Subquadratic Universal Point Sets}\label{1409}

{\small\sffamily UnsolvedMath ID 1409 · Graph Theory}

\subsection{Definitions and mathematical statement}

\paragraph{Shared notation.}Unless a section says otherwise, $\N=\{1,2,\ldots\}$,
$\Z$ is the ring of integers, $\Q,\R,\C$ are the rational, real, and complex
fields, $[N]=\{1,\ldots,\lfloor N\rfloor\}$, and $\log$ is the natural
logarithm. For real functions of a parameter tending to infinity,
$f\ll g$ and $f=O(g)$ mean $|f|\leq Cg$ eventually for some fixed $C$;
$f\gg g$ reverses this comparison; $f=o(g)$ means $f/g\to0$;
$f\sim g$ means $f/g\to1$; and $f\asymp g$ means both order bounds.
A subscript on $O$ or $\ll$ specifies allowed constant dependence.
The expression $n^{o(1)}$ in an upper bound means growth smaller than
$n^\epsilon$ for each fixed $\epsilon>0$. Local definitions govern any
exception, finite-field convention, or use of the same symbol for another
object. An asymptotic claim is not automatically an effective algorithm.



A finite set $S\subset\R^2$ is universal for $n$-vertex planar graphs if every such graph has a crossing-free straight-line drawing whose vertex positions belong to $S$. The selected $n$ points and their assignment may depend on the graph. Seek sets $S_n$ with $|S_n|=o(n^2)$. The points need not have integer coordinates. \sref{1409}{1} \sref{1409}{2}

\paragraph{Statement recorded in the supplied source.}
Do planar graphs have universal point sets of subquadratic size? \sref{1409}{1}

\paragraph{Residual task reported by the archive.}

The embedded review (2026-08-17) records: Construct an o(n\textasciicircum{}2)-size universal point set for all n-vertex planar graphs, or prove a quadratic lower bound. \sref{1409}{1}

\subsection{Short English statement}

Can one set of substantially fewer than n-squared plane points support straight-line planar drawings of every n-vertex planar graph?

\subsection{Sources}

{\small

\begin{itemize}

\item[\hypertarget{source:1409:1}{S1}] ULAM.AI, \emph{UnsolvedMath}, supplied \texttt{problems.json}, record 1409 (GRAPH-022), statement and background. The embedded literature-review date is 2026-08-17; that is the archive’s review date, not a fresh certification. Dataset landing page: \url{https://huggingface.co/datasets/ulamai/UnsolvedMath}.

\item[\hypertarget{source:1409:2}{S2}] M. Bannister, Z. Cheng, W. Devanny, and D. Eppstein, Superpatterns and Universal Point Sets, Journal of Graph Algorithms and Applications 18 (2014), 177–209. \url{https://doi.org/10.7155/jgaa.00318}. Bibliographic information imported from the supplied record.

\item[\hypertarget{source:1409:3}{S3}] M. Scheucher, H. Schrezenmaier, and R. Steiner, A Note on Universal Point Sets for Planar Graphs, Journal of Graph Algorithms and Applications 24 (2020), 247–267. \url{https://doi.org/10.7155/jgaa.00529}. Bibliographic information imported from the supplied record.

\item[\hypertarget{source:1409:4}{S4}] S. Felsner, H. Schrezenmaier, F. Schröder, and R. Steiner, Linear Size Universal Point Sets for Classes of Planar Graphs, arXiv:2303.00109 (2023). \url{https://arxiv.org/abs/2303.00109}. Bibliographic information imported from the supplied record.

\end{itemize}

}

\label{end:1409}
\end{document}
